\section{Заключение}
\label{sec:Chapter5} \index{Chapter5}

В работе исследовано влияние явной графовой структуры на способность компактных больших языковых моделей вести многошаговые цепочки рассуждений. Предложен подход, в котором корректность промежуточных шагов вывода формализуется через структурное соответствие сгенерированного графа рассуждений эталонному подграфу и используется в качестве плотного сигнала вознаграждения при дообучении с подкреплением.

\subsection{Основные результаты}
\label{subsec:conclusion-results}

В ходе работы решены все поставленные во введении задачи:
\begin{enumerate}
    \setlength{\itemsep}{\smallskipamount}
    \item \textbf{Построены датасеты с эталонными подграфами рассуждения.} Реализована предобработка RuleTaker (схлопывание вершин-правил, фильтрация по отрицаниям и NatLang-подмножеству) с публикацией датасета \texttt{KGReasoning}, а также пайплайн генерации многошаговых вопросов по графу знаний Wikidata5m (поиск $1$/$2$/$3$-hop и композиционных $3{+}2$ путей, генерация вопросов внешней LLM и валидация отсевом легко угадываемых ответов).
    \item \textbf{Формализованы метрики структурного соответствия.} Введены метрики уровня вершин и рёбер (Node/Edge Precision, Recall, $F_1$), интегральная мера Graph Edit Similarity и, для использования в обучении, рекурсивная процедура покрытия эталонного подграфа \texttt{validate\_subgraph}, явно проверяющая корректность каждого шага вывода.
    \item \textbf{Разработан фреймворк \texttt{smart-thinking-llm}.} Модульная система обеспечивает преобразование текстовых генераций в графовое представление двумя способами (прямой JSON-разбор и LLM-ассесор со Structured Output), вычисление метрик и воспроизводимую оценку.
    \item \textbf{Обоснован выбор базовой модели.} По результатам сравнительной оценки семейств Qwen3 (0.6B--32B), Llama~3.x, OLMo-2-Think и GPT-OSS для дообучения выбрана модель Qwen3-0.6B как наиболее компактная и обладающая наибольшим потенциалом улучшения.
    \item \textbf{Реализован двухэтапный пайплайн обучения.} Реализованы этап холодного старта (SFT) и этап RL с алгоритмом GRPO и композитной графовой функцией награды, развёрнутые на двух промышленных платформах~--- VERL (RuleTaker) и torchtune/\texttt{TunePPO} (Wikidata). Контрастность baseline ($\alpha=0$) и предложенного метода ($\alpha>0$) реализована единственным параметром масштабирования графовой компоненты.
    \item \textbf{Проведена экспериментальная проверка.} Пять конфигураций (база, SFT, две контрастные RL-политики и комбинированный режим SFT\,$\to$\,RL) сопоставлены на $10\,000$ примерах RuleTaker по точности и структурным метрикам, а также на пяти внешних бенчмарках.
\end{enumerate}

\textbf{Эмпирические выводы.} Гипотеза работы подтверждена. На целевой задаче RuleTaker графовая награда ($\alpha>0$) превосходит бинарный baseline ($\alpha=0$) как по точности метки ($0{,}7115$ против $0{,}6991$), так и, прежде всего, по структурным метрикам: Node~$F_1$ $0{,}2591$ против $0{,}2251$ и Edge~$F_1$ $0{,}2335$ против $0{,}1875$. Показательно, что бинарный baseline, повышая точность, \emph{не улучшает} структуру вывода (его Edge~$F_1$ не превосходит уровня базовой модели),~--- то есть структурное качество достижимо именно за счёт графового сигнала. Наилучшие результаты по всем метрикам одновременно даёт двухэтапный режим SFT\,$\to$\,RL (Accuracy $0{,}7143$, Node~$F_1$ $0{,}2671$, Edge~$F_1$ $0{,}2394$, GES $0{,}8563$). При этом дообучение на узкой логической задаче не разрушило общих способностей: на \texttt{math\_500}, \texttt{gpqa}, \texttt{mmlu\_redux} и \texttt{zebra\_logic} модели сохраняют или превосходят уровень базовой Qwen3-0.6B, а на логической головоломке \texttt{zebra\_logic} наблюдается положительный перенос навыка структурного рассуждения.

\subsection{Степень решения поставленной задачи}
\label{subsec:conclusion-degree}

Поставленная цель~--- исследование влияния графовой структуры на рассуждения LLM и разработка методов их оценки и обучения~--- достигнута в полном объёме: создан воспроизводимый инструментарий оценки графовых метрик, реализован контрастный RL-пайплайн с графовой наградой и проведено его экспериментальное сравнение с бинарным baseline. Центральная гипотеза о повышении структурного качества цепочек рассуждений под действием графовой награды подтверждена прямым сравнением по Node~$F_1$, Edge~$F_1$ и GES на $10\,000$ примерах (таблица~\ref{tab:struct-metrics}): графовая награда устойчиво превосходит baseline по всем структурным метрикам, тогда как обучение только по результату структуру не улучшает. Тем самым решены все задачи, поставленные во введении.

\subsection{Научная новизна}
\label{subsec:conclusion-novelty}

Научная новизна работы состоит в объединении трёх компонентов, ранее не использовавшихся совместно (раздел~\ref{subsec:comparison}):
\begin{enumerate}
    \setlength{\itemsep}{\smallskipamount}
    \item формализация качества цепочки рассуждений через структурное соответствие сгенерированного графа эталонному подграфу с автоматически вычисляемой мерой покрытия, не требующей ручной пошаговой разметки (в отличие от Process Reward Models);
    \item использование графовой меры соответствия непосредственно в качестве плотного сигнала вознаграждения при RL-дообучении, а не только как метрики оценки или источника фактов на инференсе (в отличие от ToG, RoG, StructGPT);
    \item обновление параметров компактной модели на основе графовой обратной связи в рамках единого фреймворка, допускающего контрастное сравнение результирующего и структурного обучения через единственный коэффициент~$\alpha$.
\end{enumerate}

\subsection{Практическая значимость}
\label{subsec:conclusion-practical}

Практическая значимость работы определяется следующим. Во-первых, разработан воспроизводимый программный инструментарий (\texttt{smart-thinking-llm} и \texttt{TunePPO}), пригодный для оценки и обучения reasoning-способностей LLM с использованием графовых метрик и развёртываемый на одной GPU. Во-вторых, предложенная методология автоматического построения плотного обучающего сигнала из структуры эталонного подграфа снижает зависимость от дорогостоящих человеческих аннотаций. В-третьих, показано, что структурно-ориентированное дообучение компактной модели не разрушает её общих способностей, что важно для практического применения малых моделей в условиях ограниченных ресурсов.

\subsection{Направления дальнейшей работы}
\label{subsec:conclusion-future}

Перспективными направлениями развития являются: проверка статистической значимости наблюдаемых различий и оценка их устойчивости по нескольким случайным инициализациям; масштабирование метода на более крупные модели семейства (1.7B--8B) для оценки зависимости эффекта от размера; полномасштабное обучение и оценка на датасете Wikidata с многошаговой графовой наградой; исследование альтернативных схем взвешивания компонент награды и более выразительных мер структурного соответствия графов; а также перенос подхода на другие задачи многошагового рассуждения за пределами RuleTaker.
